% ----------------------------------------------------
% Resumo em língua estrangeira (obrigatório no TCC)
% Preencha o título no campo {abstracttitle}.
% Use ABSTRACT, RESUMEN ou outro título conforme o idioma escolhido.
% Preencha o texto no campo {abstract}.
% Preencha os termos no campo {keywords}.
% Opcional: ajuste o rótulo em {keywordslabel}, por exemplo: Palabras clave.
% Separe os termos por ponto e vírgula.
% ----------------------------------------------------

\undfsetcontent{abstracttitle}{ABSTRACT}
%\undfsetcontent{keywordslabel}{Keywords}

\undfsetcontent{abstract}{%
	This paper presents an initial undergraduate final paper model designed to guide students in the formal organization of academic writing. The document takes the final paper itself as its theme in order to demonstrate, in an objective way, how the mandatory elements of the work are arranged, as well as the didactic division of chapters, the construction of sections, the use of citations and the presentation of references. Its purpose is to provide a simple, coherent and easily adaptable basis, allowing students to replace the illustrative passages with the actual content of their own research. By combining academic form, clear structure and pre-organized essential fields, the model seeks to facilitate the beginning of writing and to help students understand the final paper as an exercise in research, writing and scientific communication. The example also highlights where students should adapt metadata, abstracts, keywords, chapters and references while preserving the required academic structure.
}

\undfsetcontent{keywords}{%
	undergraduate final paper; academic writing; standards; UnDF}
